Reinforcement learning pricing algorithms can sustain prices above the competitive level without communicating, creating a practical problem for competition policy built around explicit agreement. This thesis asks whether real-world information and timing frictions make that concern smaller or larger. I add three frictions to the \citet{calvano2020} Q-learning pricing baseline: observation latency, noisy monitoring, and asynchronous price setting. Collusion is measured by the profit-gain index $\Delta$, where $\Delta=0$ corresponds to competitive Nash profits and $\Delta=1$ to joint-monopoly profits. Varying each friction separately produces three distinct responses in $\Delta$: a sharp initial drop for latency, a smooth decline for noise, and a step then plateau for asynchrony. A combined-friction design, replicated at $N\in\{250,500,1000\}$ sessions, then finds two interaction patterns that survive replication. Latency and noise compound super-multiplicatively, with joint disruption stronger than independent dampening would predict. Asynchrony partially rescues collusion when imposed jointly with the information frictions, reversing the sign of its standalone effect. The policy lesson is that single-friction evidence can understate the algorithmic-collusion concern, because realistic combinations of frictions may offset disruption rather than simply add to it.
